\documentclass[UTF8,fontset=none,11pt,a4paper]{ctexart}
\usepackage[left=21mm,right=21mm,top=22mm,bottom=22mm,headheight=15pt,headsep=8mm]{geometry}
\usepackage{fontspec}
\setmainfont{texgyretermes-regular.otf}[BoldFont=texgyretermes-bold.otf,ItalicFont=texgyretermes-italic.otf]
\setCJKmainfont{Songti.ttc}[Path=/System/Library/Fonts/Supplemental/,FontIndex=6,AutoFakeBold=2]
\setCJKsansfont{Songti.ttc}[Path=/System/Library/Fonts/Supplemental/,FontIndex=6,AutoFakeBold=2]
\usepackage[table]{xcolor}
\definecolor{TableHead}{HTML}{EEF0ED}
\definecolor{Muted}{HTML}{666C68}
\usepackage{booktabs,longtable,array,fancyhdr,needspace}
\usepackage{xurl}
\usepackage[unicode,hidelinks]{hyperref}
\urlstyle{same}
\linespread{1.32}
\setlength{\parindent}{2em}
\setlength{\parskip}{3pt}
\setlength{\LTpre}{8pt}\setlength{\LTpost}{10pt}
\ctexset{section={format=\large\bfseries,beforeskip=18pt,afterskip=8pt},subsection={format=\normalsize\bfseries,beforeskip=12pt,afterskip=6pt}}
\pagestyle{fancy}\fancyhf{}
\renewcommand{\headrulewidth}{.3pt}
\renewcommand{\footrulewidth}{0pt}
\fancyhead[L]{\small\color{Muted}现金管理存款归集披露及政府货币基金说明书}
\fancyhead[R]{\small\color{Muted}20}
\fancyfoot[C]{\small\color{Muted}\thepage}
\hypersetup{pdftitle={现金管理存款归集披露及政府货币基金说明书},pdfauthor={金融决策模拟材料}}
\begin{document}
{\noindent\LARGE\bfseries 现金管理存款归集披露及政府货币基金说明书\par}
\vspace{7pt}{\noindent\small\color{Muted}中文模拟参考文件\par}\vspace{10pt}
\section*{第一部分　银行存款归集项目披露}
\subsection*{定义}
本披露中的“平台”“我们”按语境包括经纪服务公司、清算服务公司及关联公司。“你”或“账户所有人”指开户申请列明的所有人。存入账户、收到而待投资、待提款或待进一步指示的款项称为“现金余额”。


项目银行存款是现金余额已经归集至参与银行形成的存款；货币市场基金溢出部分不属于项目银行存款。适用于不同现金管理、退休及健康储蓄账户的银行清单可能不同。本文件的本次账户记录涉及普通现金管理账户。


\subsection*{概述}
选择存款归集项目后，账户现金余额自动存入或从一个或多个参与银行的计息存款账户取出。参与银行不可接受存款，或现金余额超过项目可用容量及限额时，有关现金归集至货币市场基金溢出安排。


是否选择参加项目，由开户申请或现有账户变更指示决定。平台不替你决定是否采用该核心头寸，也不替你接受银行清单及其变更。你可以随时将现金余额移出项目，或停止采用项目，而不承担项目退出罚金。


无法存入银行，或平台判断可能不能适用存款保险的余额，按下文规则放入货币基金溢出安排。这会改变有关余额的保险适用、收益计算和存取机制。


归集安排不作为长期投资选择。账户有其他可用核心头寸，具体可用性按账户协议处理。


\subsection*{项目运作——存款}
每个营业日，账户现金余额自动归集至平台代表你及其他参加客户在参与银行设立的项目存款账户。营业日不包括银行假日及纽约证券交易所休市日，例如耶稣受难日。


归集从账户清单中的主要参与银行开始，直至该行的项目存款达到最高存款限额。剩余现金依清单顺序进入下一个可接受存款的银行，直至没有剩余现金，或各可用银行均已达到限额。


各可用银行均达到限额后，剩余现金进入货币市场基金溢出安排。在个别有限情形下，主要银行仍可能持有超过最高存款限额的余额，超出保险限额的部分不受FDIC保障。


\subsection*{最高存款限额}
项目最高存款限额始终为当时适用的个人退休账户标准最高存款保险金额的98\%。例如标准保险金额为250,000美元时，项目最高存款限额为245,000美元。


这一项目配置限额与FDIC保险的实际合并计算是不同事项；本人的其他同银行同所有权类别存款仍须合并。


\subsection*{提款}
营业日结束时，账户需要资金支付借项，先使用货币基金溢出余额。仍需要资金，或没有溢出余额时，项目银行存款自动在下一营业日归集回账户。


某些日间实际或预计交易产生的借项，改为使用当日赎回溢出基金或当日提取银行存款的款项结算。


银行提款顺序与存入顺序相同：先从主要银行提取，再顺次从后续可用银行提取，直至借项付清或全部项目银行存款已经取回。留在各银行的存款不因该过程重新分配。因此，主要银行可以低于配置限额，而其他银行仍处于配置限额。


\subsection*{利息入账}
银行收到项目存款当天开始计息，提取当天停止计息。账户报表分别列明每个参与银行的应计利息。


银行存款每日计息，每月复利，于每月最后一个营业日记入项目存款账户。利息入账后归集至你的账户，可形成新的现金余额，再按存款规则归集。


溢出基金通常每日宣布股息、每月支付股息。形成的现金余额同样按有关归集规则处理。


\subsection*{存款保险限额变更}
适用于个人退休账户或普通单一所有权存款账户的标准保险金额变更时，平台按变更生效日重新确定有关账户的项目配置限额。


限额提高时，平台可以提高一个或多个银行的余额，减少清单中需要使用的银行数；限额降低时，可将余额分配至后续银行。能够参与的银行容量及溢出机制继续适用。


\subsection*{所有权记录}
平台保存反映各客户权益的记录。银行账户可以以平台或清算代理的名义设立，但记录须反映客户的受益所有权。FDIC保险依客户的所有权类别及适用规则计算，而不是对平台名称下所有客户合计只给一份保险限额。


\subsection*{参与银行}
各参与银行可以设定其可接受的项目存款总额，即银行存款容量。达到容量后，该银行可以不再接受新的项目存款。在容量尚未用满时，银行也可能因为其他原因无法接受存款。


\subsection*{银行状态}
可用：原则上可以接受项目存款，除非适用以下其他状态。


不可用：不能接受额外存款。该状态本身不改变已经持有的存款，但阻止新增存款。在没有可用溢出安排，且该银行为清单中唯一银行，或主要银行之外各银行均不可用、已排除或待激活等特定情形下，文件允许例外处理。


溢出银行：暂不能担任主要银行，已经存入的存款不因此变化。这里的银行状态“溢出”与货币市场基金溢出是两个不同概念。该银行是有关主清单唯一银行时，可以适用例外。


待激活：已经加入主清单，但尚不能接受项目存款。


客户已排除：你决定不采用该银行。银行仍列在你的清单中，但不接受本账户新归集存款；这只影响你的账户，不改变该银行在整个项目中的参与状态。排除须联系平台代表。


\subsection*{账户银行清单}
平台每天根据有关账户类型的银行主清单生成账户银行清单。不同开户日期，或不同账户类型，可能对应不同清单。


清单中的次序决定银行使用次序。第一家为主要银行。平台保留提前通知后修改清单的权利；分配初始清单后继续使用账户，或收到变更通知后继续使用，表示接受清单。


账户建立后，可以按适用协议排除一家或多家银行；存在其他清单时，也可以申请采用另一清单。


\subsection*{新账户}
新开户人可以按开户申请及客户协议选择归集。平台可因任何原因暂停或恢复向新账户提供该项目；暂停时可以选择其他可用核心头寸。


同一天新建立的同类账户采用当日生成的初始清单。清单由有关主清单中的部分或全部银行组成，并依主清单变更自动更新。


\subsection*{现有账户转入归集项目}
符合资格的现有账户可以按相关客户协议，在可用货币基金核心头寸与银行归集之间变更。决定转入归集时，采用当日分配给新账户的清单，并可依文件规定排除银行。


\subsection*{已采用项目的账户}
登入账户的持仓页面，可以查看当前清单、每行项目存款及银行状态。有银行状态、清单或项目相关问题时，可联系平台代表。


\subsection*{客户发起清单变更}
参加项目至少须保留一家可用银行。排除某银行使其在本账户中不再接受存款，但仍显示其状态。该排除指示可以撤回，只要银行仍参加项目且可以接受存款。


如果申请换清单，原已排除银行的指示原则上延续。但不能排除新清单的主要银行；要维持该排除，须使用主要银行不同的另一清单。


排除仍有余额的银行时，该余额被视为账户现金，并按规则重新分配至其他银行或溢出基金。改用另一清单时，原全部项目存款也按新清单重新分配。


\subsection*{银行加入及退出}
新银行可以加入清单，原银行可以退出。单纯增加或删除银行不重新排列其余银行的顺序。退出后，该银行不再接受新存款；平台提供更新清单，并将原存款转到其余银行或溢出安排。


主要银行退出时，从清单中指定另一可用银行为主要银行。每个参与账户至少须有一家可用银行；银行退出与客户排除指示共同导致无可用银行时，现金进入溢出基金。


\subsection*{银行状态变更}
新增银行可能先显示待激活，通常在加入清单约三十日后才能接受存款；客户可在待激活期间排除该银行。


银行由待激活、不可用或客户已排除状态恢复可用后，主要银行超过最高配置限额的余额可自动转往恢复可用的银行，按存款次序处理。


\subsection*{其他项目变更}
条件允许时，平台提前通知影响账户的变更。银行不愿或不能接受存款、项目参与终止，或银行持续经营及保险资格出现疑问等特殊情况下，你指示平台可以不事先通知，在日间修改清单、限制归集金额，或将余额移至其他银行或溢出基金。


没有可用银行时，部分或全部现金可以进入溢出基金。收到变更通知后继续使用账户，表示同意变更。


平台也可能必须把核心头寸改为另一可用核心头寸。提前书面通知后，你未在规定期间内提出不同指示，视为同意并指示平台提取银行及溢出余额、转入新核心头寸。


新核心头寸可能不适用FDIC保险，而可能适用SIPC保护。平台不保证新头寸提供相同或更高收益；实际收益可能较低，平台所得报酬也可能不同或更高。


签署开户申请或指示采用归集，表示已经收到及阅读本披露，并同意必要时变更核心头寸。转入独立第三方顾问或平台关联顾问的有关服务时，可以按文件终止参与并清算有关余额。


你同意平台向项目服务方提供姓名、社会安全号码或税务号码、地址及出生日期等资料。银行等服务方仅为履行其本身或平台的法定及监管义务而使用这些资料。


\subsection*{现金余额的访问}
只能通过本平台账户访问现金，不能直接联系参与银行提取你的项目存款。送达平台的征税、扣押等法律程序可以作用于有关存款。


银行存款是银行对你的直接债务，不是平台的债务。平台不保证参与银行的财务状况。联邦银行规则允许银行在提款前要求七日通知。你在项目存款账户的权益不能转让。


\subsection*{存款账户}
参与银行采用活期存款账户持有项目存款。你在银行没有独立活期子账户，存款与其他项目客户存款合并持有。文件为便于描述而将分配给你的部分称为“你的”活期子账户。


\subsection*{报表及确认}
每期报表显示核心头寸余额、各银行在报表期最后营业日的存款及溢出基金余额，列明归集进出交易，并分别列明银行及溢出收益。报表替代每笔归集的单独确认。


你负责监测每个银行全部同类别存款，包括通过其他安排存入的款项。银行清单变化可能影响实际保险覆盖。


\subsection*{货币基金溢出}
各可用银行达到配置限额，或项目及银行没有足够存款容量时，不能存入的现金进入溢出基金。


账户借项先使用溢出基金，再从主要银行及后续银行提款。某些日间实际或预计交易使用当日赎回及提款款项；其他借项按下一营业日提款安排结算。


银行恢复容量之后，已在溢出基金中的款项不会因此自动转回或重新平衡到银行。除用于支付借项、提款或适用再平衡事件外，仍留在溢出基金。


持仓页面显示溢出余额；报表显示相关余额、归集记录及收益。溢出部分不适用FDIC存款保险，可以在适用限额内受SIPC保护。适用时提供有关溢出基金说明书。


\subsection*{银行存款利率}
利率可以随时改变，并受到银行支付利率、平台费用、账户类型及文件列明的其他因素影响。项目利率可能低于直接在参与银行开户的利率，或其他非FDIC投资及货币基金收益。


不同账户类型可采用不同利率；本人其他账户取得的利率不必与本账户一致。具体当时利率可向平台或利率页面查询。


\subsection*{溢出基金收益}
溢出基金通常采用七日年化收益指标，随基金投资收入及其他因素变化。有效收益率还考虑年度复利。


货币基金取得利息、股息及其他收入，扣除费用后向持有人分配；实现资本收益扣除损失后，也可以分配。通常每日宣布、每月支付股息。收到资金日开始累计股息，提取日停止。


\subsection*{FDIC保险及SIPC保护}
项目存款须与本人在同一银行、同一所有权类别中的其他存款合并计算，包括其他归集项目、储蓄、支票、货币市场存款及银行直接发行存单。


单一所有权账户通常最高250,000美元；共同账户每位共同所有人份额通常最高250,000美元。退休账户和健康储蓄账户依各自所有权分类处理；信托存款有另外规则。不同所有权类别依法可以分别保障。


平台不掌握本人全部非项目存款，也不负责替本人监测是否超过保险限额。本人须监测合并金额。超过适用限额的资金不受FDIC保险。


保险金额包括符合规则的本金及应计未付利息。银行倒闭需要保险支付时，由FDIC提供；没有特定必须向本人提供款项的时间期限，资金访问可能显著延迟，也可能要求本人向FDIC及平台提交文件。


账户证券及溢出基金是投资产品，不适用FDIC保险，没有银行或政府本金保证，可能损失本金。SIPC对客户资产保护的适用上限为500,000美元，其中现金请求上限250,000美元；不赔偿投资市场价值波动。


现金只有在成为参与银行持有的项目存款之后才适用FDIC。在平台持有或银行进出途中，一般按SIPC规则处理。另选货币基金或其他核心投资头寸时，也不因此成为FDIC存款。


\subsection*{平台与其他机构的关系}
平台为经纪账户提供相关服务；清算服务公司作为代理在银行设立合并项目账户，并办理归集。项目管理方提供记录及技术服务。退休或健康储蓄账户如适用，另有托管主体。


平台可以同时是项目管理方或银行的客户，并存在其他金融往来；项目管理方也可以是参与银行。以后关联公司可以被指定为项目管理方、参与银行或两者。


平台母公司对某参与银行持有少数、非控制性权益。关联基金管理公司负责管理溢出基金，并取得管理及其他费用。


\subsection*{平台与其他机构取得的利益}
银行向平台支付项目费用，通常按每日项目存款合计余额计算。费用因银行、账户类型及时间而异，存款增加通常使费用增加；同一银行在其他类似项目支付的费用也可能不同。


费用可达到该银行有关项目存款余额的年化4\%。平台可降低或免除部分费用；超过4\%的提高会另行通知。该费用由参与银行支付给平台。


银行付给平台的费用，按银行合同“费率”乘以存款余额，再减银行向项目账户支付的全部利息计算。合同费率通常与有效联邦基金利率或目标利率有关，可能按余额分档，并处于相关联邦基金利率上下0.75个百分点的范围内。


项目管理方为技术、会计及监管协助服务收取平台支付的年费，按每日项目存款合计余额的0.010\%计算，并设固定上限。管理方同时是参与银行时，银行付给平台的费用与平台付给管理方的费用各自独立。


银行将项目存款用于现有及新增贷款和投资。其收入来自这些资产收入与支付给存款人利息及平台费用之间的差额。银行可能支付低于市场或直接开户存款的利率。平台还可能因关联持股取得股息或利润。


溢出基金管理机构取得基金管理和运营费用，部分费用可以与平台分享。


\section*{第二部分　政府货币市场基金招募说明书}
政府货币市场基金　SPAXX


招募说明书资料日期：2026年6月29日。证券监管机构未批准或否决本证券，亦未认定说明书准确完整；相反陈述属犯罪行为。


\subsection*{投资目标}
在兼顾资本保存及流动性的条件下，寻求尽可能高的当期收入。


\subsection*{费用表}
购买、持有及出售可能产生下列费用。经纪佣金及向中介支付的其他费用可能未列入。


\begingroup\small\setlength{\tabcolsep}{4pt}\renewcommand{\arraystretch}{1.35}
\begin{longtable}{>{\raggedright\arraybackslash}p{0.68\dimexpr\linewidth-4\tabcolsep\relax}>{\raggedright\arraybackslash}p{0.32\dimexpr\linewidth-4\tabcolsep\relax}}
\toprule
\rowcolor{TableHead} 年度运营费用 & 占投资价值比例 \\
\midrule\endfirsthead
\toprule
\rowcolor{TableHead} 年度运营费用 & 占投资价值比例 \\
\midrule\endhead
\midrule\endfoot
\bottomrule\endlastfoot
管理费 & 0.25\% \\ \addlinespace[3pt]
分销／服务12b-1费 & 无 \\ \addlinespace[3pt]
其他费用 & 0.17\% \\ \addlinespace[3pt]
总年度运营费 & 0.42\% \\
\end{longtable}\endgroup
股东直接费用：无。


费用示例：假设份额年收益5\%，费用保持表中水平，每投入10,000美元并在期末全部出售，估计费用分别为一年43美元、三年135美元、五年235美元、十年530美元。例子用于比较费用，不代表实际或预计收益和费用。


\subsection*{主要投资策略}
通常至少99.5\%总资产为现金、美国政府证券或完全以现金／政府证券抵押的回购协议。包括国会特许或支持、但并非美国财政部发行或保证的机构证券。


遵守货币市场基金投资质量、期限、流动性及分散的监管标准。此外，通常至少80\%资产投资美国政府证券及该等证券回购协议。


政府证券可能由财政部或政府机构发行或保证，其支持可以是财政部完整信用、向财政部借款权，或发行机构本身。房利美、房地美及联邦住房贷款银行等证券并非财政部发行或保证。证券可付固定、可变或浮动利息，通常到期偿还本金。


回购协议指按一个价格买证券，同时同意按约定价格卖回，可用政府证券或现金抵押。


\subsection*{主要风险}
利率增加可能使货币市场证券价格下降；低或负利率可能影响收益；发行人或提供保证、缩短期限结构的主体信用下降，可能降低价格。部分政府机构及政府支持企业证券不由美国政府保证，也没有美国完整信用支持。


投资可能亏损。基金以每份1美元为目标，不能保证达到。基金不是银行账户，不受FDIC或其他政府保险。提供者及关联方没有补偿损失的义务，不应预期其在市场压力期间提供支持。


本政府基金不在出售份额时收取赎回费用。


\subsection*{业绩}
以下显示每年的变化。过往表现不表示未来表现。当前资料可向基金机构查询。


\begingroup\small\setlength{\tabcolsep}{4pt}\renewcommand{\arraystretch}{1.35}
\begin{longtable}{>{\raggedright\arraybackslash}p{0.68\dimexpr\linewidth-4\tabcolsep\relax}>{\raggedright\arraybackslash}p{0.32\dimexpr\linewidth-4\tabcolsep\relax}}
\toprule
\rowcolor{TableHead} 年份 & 年度收益 \\
\midrule\endfirsthead
\toprule
\rowcolor{TableHead} 年份 & 年度收益 \\
\midrule\endhead
\midrule\endfoot
\bottomrule\endlastfoot
2016 & 0.04\% \\ \addlinespace[3pt]
2017 & 0.51\% \\ \addlinespace[3pt]
2018 & 1.47\% \\ \addlinespace[3pt]
2019 & 1.84\% \\ \addlinespace[3pt]
2020 & 0.26\% \\ \addlinespace[3pt]
2021 & 0.01\% \\ \addlinespace[3pt]
2022 & 1.31\% \\ \addlinespace[3pt]
2023 & 4.78\% \\ \addlinespace[3pt]
2024 & 4.92\% \\ \addlinespace[3pt]
2025 & 3.95\% \\
\end{longtable}\endgroup
显示期间最高季度为2023年12月31日结束季度，1.26\%；最低为2022年3月31日结束季度，0.00\%；2026年3月31日为止年初至今0.82\%。


截至2025年12月31日，一年、五年、十年平均年度收益为3.95\%、2.97\%、1.89\%。


\subsection*{顾问、交易及税务}
基金经理为Fidelity Management \& Research Company LLC，其他顾问提供次级顾问服务。


可通过经纪、共同基金、退休账户或投资专业人员买卖，方式包括互联网、电话和邮寄。每份买卖价格是适当形式订单收到后下一次计算的净值。没有购买最低金额。


纽约证券交易所开市日营业。交易所关闭，但纽约联储及组合主要市场营业，并认为市场足以满足请求时，也可能营业。


分配一般适用联邦所得税及可能的州地方税，通常为普通收入或资本收益。递延税退休账户可能提款时才计税。


基金、顾问、分销商及关联方可能向银行、经纪商、退休计划及服务机构支付销售与服务报酬，可能使中介有推荐本基金的激励。


\subsection*{投资详情}
顾问强调维持稳定1美元价格、流动性及收入，并按质量、期限、流动性和分散标准交易。货币市场证券是优质短期债务，可以包括固定、可变或浮动利率证券。


公共卫生、军事冲突、恐怖主义、政府限制、政治或环境灾害可能影响表现。收益每日随利率及市场改变；利率大涨或发行人信用降低可能影响净值。美国政府不保证净值或收益。


期限较长或某些证券更敏感；短期证券主要对短期利率变化反应。利率上升可能使发行人更难及时偿本付息。短期资产到期后以不同利率再投资，基金收益也改变。


顾问可能通过费用减免或报销避免负收益，但不保证。发行人、交易对手、保证者或附带赎回结构的状况改变，可能导致违约或结构失效。部分政府义务未有国家完整信用保证。


为防御市场、经济、政治等情况，基金可能临时采用其他策略，包括较多资产未投资；不产生收入的资产可能降低收益，并影响目标达成。


\subsection*{基本政策与份额估值}
投资目标是基本政策，须股东批准才能改变。


基金按摊余成本方法估值，取购入成本并按应计利息和折溢价调整，旨在保持每份1美元；该方法不保证价格始终稳定。净资产除流通份额数为每份净值。基金按自己的营业日及公布时点计算。


\subsection*{买卖份额的一般事项}
货币基金设计为提供流动现金选择，预期经常买卖，董事会未制定一般抑制频繁交易的政策。但基金可拒绝、限制或取消被认为干扰管理或不符合基金利益的购买及转换，并可设更严格条件。


份额一般仅向美国居民提供，可能要求证明资格。没有本说明书份额的最低余额或购买额，不收销售费。


\subsection*{买入}
适当形式订单收到后，以之后下一次净值成交。获授权中介收到正确订单，视为基金收到。基金可暂时或永久停止发行或限制购买，包括暂停赎回期间。


未收到并收妥付款，可以取消购买，你可能承担损失或费用。电汇订单表示当日到账时，资金须在联邦储备电汇系统关闭前按要求送达指定银行。反洗钱等规则可使订单暂停、限制或取消，并保留款项。


\subsection*{卖出}
适当形式赎回订单收到后，以之后下一次净值成交。通常下一营业日处理；立即付款不利于基金时，支付款项最多可延至七日。中介负责及时转交订单。


某些申请须书面及签名担保，包括出售超过100,000美元、记录地址十五日内改变或支票寄不同地址、支付给非持有人，以及转往不同登记的账户。银行、经纪、信用合作社等可提供担保；公证员不能替代。


原购买款未收妥，赎回可能延迟；交易所非假日关闭、交易受限或SEC允许时，也可暂停或延期。顾问认为符合基金利益时，可能用证券或其他财产支付。


可用支票赎回，但支票超过持仓价值可能被退回及收费。未兑现赎回支票不产生利息；六个月未兑现时，可能按当日下一次净值再投资。


反洗钱规则可能暂停、限制、取消或处理后扣留赎回款。董事会在1940年投资公司法Rule 22e-3规定情形下，可决定永久停止赎回并清算。


\subsection*{赎回支付政策}
通常下一营业日或适用同日结算支付。自动卖出本基金以买另一基金时同步结算。法规允许的特殊延迟最多七日。


通常利用现金、证券销售款及其他流动资产支付；正常和市场压力中都可以出售组合证券、采用实物赎回或允许的短期借款等方式。


\subsection*{份额转换}
按资格可把份额换为其他基金；转换按各基金的购买销售规则及下一次净值处理。其他基金可能有不同目标、风险及费用，须阅读其说明书。基金保留修改或终止转换功能的权利。


\subsection*{账户功能与政策}
可提供自动投资、自动提款、支票及其他服务，依账户资格及平台规则。变更资料、交易指示及转移按要求提交。身份及税务资料缺失可能限制交易。


\subsection*{股息及资本收益分配}
基金每日宣告、每月支付收入股息；资本收益按需要分配。可选择现金或再投资，未另选时通常再投资。收到现金支票未投递时，可能改为再投资；未兑现支票不计息，六个月后可能投入新份额。


\subsection*{税务后果}
非退休账户的分配可能产生联邦及州地方税。来自某些政府证券、并符合条件的部分股息，可能豁免州地方税。


股息及短期资本收益通常是普通收入，长期资本收益分配通常按资本收益。基金主要收入来自利息，股息通常不符合个人长期资本收益优惠税率。


现金或再投资都不改变当期通常应税。某些十二月分配一月收到，税务上按十二月三十一日收到。


\subsection*{基金管理}
共同基金汇集股东资金，按指定目标投资。管理人负责投资和业务，地址245 Summer Street, Boston, Massachusetts 02210。截至2025年底，管理人全权资产约5.6万亿美元，连同关联方约7.1万亿美元。


英国、香港及日本关联机构提供有关发行人研究及次级顾问服务。员工发表的市场观点仅为当时个人意见，可改变，不表示机构投资建议或基金交易意图。


基金每月支付按平均净资产计算的年度0.25\%管理费。管理人向次级顾问付费；基金董事会批准有关合同的依据见相应定期报告。顾问可能暂时减免费用并保留符合条件的偿还权，减免可降低费用、提高表现。


\subsection*{分销}
基金不同份额类别共用目标和组合。分销商提供销售，中介可因记录保存、行政、退休计划及销售服务取得报酬。


12b-1计划允许管理人从管理费收入、过去利润或其他资源支付销售及股东支持费用。若被认定持续从份额资产支付，会增加投资成本，可能比其他销售收费更高。


分销商可能提供商品、折扣、券等促销，资格见具体条件。没有销售人员或其他人获授权提供说明书及补充资料以外的陈述；不得将其他陈述视为基金授权。违法对象或地区不构成本文件要约。


\subsection*{财务摘要}
以下资料用于说明基金份额过去五个会计年度的财务历史。部分项目以每一份基金份额表示；总回报假定全部股息及其他分配再投资。年度数据经独立注册会计师事务所审计，其报告及基金财务报表载于年报，年报可免费索取。


\begingroup\small\setlength{\tabcolsep}{4pt}\renewcommand{\arraystretch}{1.35}
\begin{longtable}{>{\raggedright\arraybackslash}p{0.3\dimexpr\linewidth-12\tabcolsep\relax}>{\raggedright\arraybackslash}p{0.14\dimexpr\linewidth-12\tabcolsep\relax}>{\raggedright\arraybackslash}p{0.14\dimexpr\linewidth-12\tabcolsep\relax}>{\raggedright\arraybackslash}p{0.14\dimexpr\linewidth-12\tabcolsep\relax}>{\raggedright\arraybackslash}p{0.14\dimexpr\linewidth-12\tabcolsep\relax}>{\raggedright\arraybackslash}p{0.14\dimexpr\linewidth-12\tabcolsep\relax}}
\toprule
\rowcolor{TableHead} 截至4月30日年度 & 2026 & 2025 & 2024 & 2023 & 2022 \\
\midrule\endfirsthead
\toprule
\rowcolor{TableHead} 截至4月30日年度 & 2026 & 2025 & 2024 & 2023 & 2022 \\
\midrule\endhead
\midrule\endfoot
\bottomrule\endlastfoot
期初每份净值（美元） & 1.00 & 1.00 & 1.00 & 1.00 & 1.00 \\ \addlinespace[3pt]
每份净投资收入（美元） & 0.037 & 0.045 & 0.049 & 0.027 & 少于0.0005 \\ \addlinespace[3pt]
每份已实现及未实现损益 & 少于0.0005 & 少于0.0005 & 少于0.0005 & 少于0.0005 & 少于0.0005 \\ \addlinespace[3pt]
来自运营的每份合计（美元） & 0.037 & 0.045 & 0.049 & 0.027 & 少于0.0005 \\ \addlinespace[3pt]
净投资收入分配（美元） & (0.037) & (0.045) & (0.049) & (0.027) & 少于0.0005 \\ \addlinespace[3pt]
期末每份净值（美元） & 1.00 & 1.00 & 1.00 & 1.00 & 1.00 \\ \addlinespace[3pt]
总回报 & 3.72\% & 4.57\% & 5.05\% & 2.72\% & 0.01\% \\ \addlinespace[3pt]
减免前费用率 & 0.42\% & 0.42\% & 0.42\% & 0.42\% & 0.42\% \\ \addlinespace[3pt]
扣除费用豁免后的费用率 & 0.42\% & 0.42\% & 0.42\% & 0.42\% & 0.10\% \\ \addlinespace[3pt]
扣除全部减免后的费用率 & 0.42\% & 0.42\% & 0.42\% & 0.42\% & 0.10\% \\ \addlinespace[3pt]
净投资收入率 & 3.64\% & 4.45\% & 4.94\% & 2.68\% & 0.01\% \\ \addlinespace[3pt]
期末净资产（千美元） & 420,114,025 & 366,312,998 & 297,752,603 & 236,784,703 & 230,027,515 \\
\end{longtable}\endgroup
A：净投资收入受所持基金或ETF宣布股息时点影响，所持基金自身的净投资收入率不计入本基金的该项比率。


B：有关金额少于每份0.0005美元。


C：有关期间没有部分费用减免时，总回报会较低。


D：所持基金或ETF的费用未计入本基金费用率，本基金间接承担其相应份额。


E：减免前费用率不反映补偿、豁免或其他抵扣安排，不代表发生减免期间实际承担的金额。


\section*{本次账户资料及办理记录}
\subsection*{A：银行存款归集}
\textbf{转入资金的去向：} 平台按存款归集规则转入参与银行，形成银行存款


\textbf{开始比较使用的收益指标：} 存款年化收益率APY为1.84\%


\textbf{指标的含义：} 按所列存款利率及复利规则计算的年化收益率，存款利率可调整


\textbf{基金内部运营费：} 不购买该基金，无该项基金费


\textbf{资金保障机制：} 符合记录及存款归集条件时，按参与银行的FDIC存款保险规则保障


\textbf{资金运作：} 参与银行依存款规则持有及计息


\textbf{价值及账户付款：} 从银行存款中提取相应金额处理账户付款


\textbf{使用账户的付款工具：} 借记卡、支票、账单支付、转账


\subsection*{B：政府货币市场基金}
\textbf{转入资金的去向：} 自动买入政府货币市场基金份额


\textbf{开始比较使用的收益指标：} 基金七日年化收益率为3.28\%


\textbf{指标的含义：} 根据最近7日净收益折算的年化指标，已反映基金运营费，随实际基金收益变化


\textbf{基金内部运营费：} 总年度运营费率0.42\%，从基金资产承担，不另外按账户余额开账单


\textbf{资金保障机制：} 基金份额不是银行存款，不属于FDIC存款保险


\textbf{资金运作：} 主要投资现金、美国政府证券及符合规则的回购协议


\textbf{价值及账户付款：} 基金以每份1美元的稳定价格为目标，账户付款时自动赎回相应份额；该价格目标不是本金保证


\textbf{使用账户的付款工具：} 相同，使用自动赎回功能处理付款


这里的两个收益指标是本次场景开始时采用的起始资料，不是固定12个月的合同利率。B列示的七日年化指标与A的APY采用不同计算口径；之后的利率、基金实际收益和余额随各自机制及你的支出变化。


\textbf{公开历史资料。} 你查阅的现金产品资料附有以下两项历史记录。它们涉及的机构和产品均不是本次账户的参与银行或政府货币基金：


\subsection*{资产及事件}
\textbf{2008年9月：} 货币基金R持有企业发行的债务证券，其中包括雷曼证券；基金每份净值降至0.97美元


\textbf{2023年3月：} 银行V倒闭，存款账户由接管及处置安排处理


\subsection*{当时的处理记录}
\textbf{2008年9月：} 基金宣布暂停支付赎回款，随后进入相关处置安排


\textbf{2023年3月：} 主管部门采用系统性风险例外保护该行全部存款，存款人自3月13日起可以使用资金；该事件包括超过常规保险限额的存款


\subsection*{办理确认事项}
A的FDIC保险按每个参与银行、每位存款人和账户所有权类别计算，通常每类最高250,000美元。该银行其他同类别存款需合并计算；本次已核对合并余额在额度内，平台保留满足归集保障条件的记录。保险对应参与银行倒闭时的合格存款，不是给基金投资提供的保障。


B采用政府货币市场基金规则，至少99.5\%的总资产配置于现金、政府证券或以现金／政府证券充分抵押的回购协议。基金收入随短期资产收益变化，分配的收益留在账户并按核心头寸机制继续运作。提款所用份额按照实际适用的赎回价值结算。


两项日常支付工具相同，本次没有因选择某项而增加专属客服或理财服务。你选择的是现金的持有和收益机制；以后可以按平台规则更换核心头寸，但本次不把将来换项安排预定为结果。

\end{document}
